% !TeX root = ../main.tex
\chapter{Introduction}
\section{Background of the Study}
A functional legal system requires strict consistency so that citizens can clearly understand their rights and obligations \cite{fuller1964the}. However, the sheer volume of modern legislation has grown so rapidly that human practitioners can no longer track every active rule \cite{ruhl2015measuring}. As governing bodies continuously pass new regulations to address emerging issues, they frequently enact laws that inadvertently contradict older ones \cite{katz2024gpt4}. Because the total number of documents is overwhelming, identifying these hidden conflicts through manual review is computationally impossible for legal staff. This issue is particularly severe in local government units, where historical ordinances are often scattered across fragmented files rather than maintained in a structured digital database \cite{dias2022state}.

Within the Philippine legal framework, local governments like Davao City operate under delegated authority and can only enact laws permitted by the national government \cite{philgov1991ra7160}. Due to this structural limitation, the Magtajas v. Pryce Doctrine strictly prohibits any local municipal ordinance from contradicting established national statutes \cite{scp1994magtajasvpryce}. Consequently, in the local context, city councils must rigorously audit every newly proposed draft ordinance to ensure it logically aligns with existing laws before final approval. If reviewers overlook a normative conflict, such as a local ordinance prohibiting an activity that a national law explicitly permits, the resulting local legislation is legally void. Passing invalid laws creates severe jurisdictional disputes, wastes administrative resources, and ultimately compromises the constitutional rights of the public.

To address these legal risks, the legislative process must utilize an ex-ante evaluation: a term derived from Latin meaning 'before the event.' Unlike an 'ex-post' review, which happens after a law is already in effect and often after legal damage has occurred, an ex-ante audit acts as a first line of defense. By checking a proposed ordinance for conflicts while it is still in the drafting phase, the Sangguniang Panlungsod can proactively align local rules with national statutes. This preventative approach is consistent with modern Philippine standards like Republic Act No. 11032 or the Ease of Doing Business and Efficient Government Service Delivery Act of 2018, which encourages government agencies to perform regulatory assessments to ensure that new rules do not create unnecessary or conflicting burdens on the public.

The decentralized drafting process of local ordinances further complicates manual review because the resulting texts are often highly unstructured \cite{spdc2022ra00122}. City councils frequently splice and amend different drafts, producing verbose documents that blend formal legal English with regional dialects like Cebuano and administrative "Taglish" \cite{dreisbach2021language}. Furthermore, because municipal archives heavily rely on digitized PDF scans, the text suffers from severe Optical Character Recognition (OCR) errors. While OCR is typically 99 percent accurate for clean English text, studies show accuracy can drop to 78 percent for low-resource languages, with Cebuano specifically seeing word error rates around 6.9 percent \cite{oard2003desperately}. These errors cause an ``error cascade'' that shatters standard vocabulary, making it difficult for traditional search engines to semantically differentiate between a statutory obligation and a strict prohibition.

In light of these challenges, the proponents seek to develop a specialized application capable of retrieving relevant legal documents and comparing them against proposed ordinances to detect potential conflicts.


\section{Problem Statement}
This study is conducted to investigate the use of a multi-stage architecture to automate the manual auditing of proposed local legislation in Davao City. 

Specifically, the study aims to answer the following research questions:

\begin{itemize}
    \item What are the specific linguistic parameters and structural markers that define a "semantic conflict" between a proposed local ordinance and existing local or national laws?
    \item How can a Coarse-to-Fine system be designed to efficiently retrieve relevant national and local statutes before performing deep logical inference on constrained LGU hardware?
    \item Which combination of a hybrid Information Retrieval (IR) algorithms and Natural Language Inference (NLI) models provides the highest accuracy when processing the specific hierarchy of Philippine legal texts?
    \item What is the empirical performance of the proposed pipeline when evaluated against a manually curated Ground Truth dataset of local and national conflict pairs?
    \item How usable and viable is the final system interface for key stakeholders within the Davao City Sangguniang Panlungsod?
\end{itemize}

\section{Objectives of the Study}
The primary goal of this research is to develop a computationally efficient, two-stage semantic conflict detection system for proposed Davao City ordinances. 

Specifically, this study aims to:
\begin{itemize}
    \item Identify the specific linguistic characteristics and structural constraints of Davao City’s legislative archive that contribute to semantic ambiguity and retrieval failure;
    \item Design a Coarse-to-Fine computational pipeline that decouples document retrieval from logical inference to optimize resource usage for local deployment;
    \item Implement a hybrid Information Retrieval (IR) algorithm and Natural Language Inference (NLI) system specifically calibrated for the Philippine legal hierarchy;
    \item Evaluate the system’s reliability by comparing its automated detections against a verified set of known legal conflicts to measure how accurately it identifies contradictions; and
    \item Develop a user-friendly web interface for the Davao City Sangguniang Panlungsod designed to accept uploaded draft documents and provide real-time conflict alerts.
\end{itemize}


\section{Significance of the Study}
The Davao City Sangguniang Panlungsod stands as the primary beneficiary of this automated consistency checking system. By transitioning from manual review to computational auditing, the tool saves legislative staff hundreds of hours of reading time. This technological intervention actively prevents the city from enacting logically flawed laws that invite future litigation or administrative deadlocks. It empowers the local government to draft regulatory measures with a significantly higher degree of legal certainty and operational speed. Ultimately, the system modernizes the legislative lifecycle by embedding rigorous logical checks directly into the initial drafting phase.

This study also provides substantial value to legal researchers and the local Natural Language Processing community. It establishes an empirical baseline for processing messy, localized Philippine legal corpora that frequently contain "Taglish" and archaic Spanish loanwords. The research demonstrates that efficient, edge-deployable encoder models can successfully parse OCR-degraded documents without relying on expensive, cloud-based supercomputers. Beyond academic and governmental applications, the general public benefits directly through the consistent implementation of coherent public policy. The tool acts as a procedural safety net, protecting citizens from the legal traps and frustrations of contradictory municipal obligations.

\section{Scope and Limitations}
The scope of this study focuses on developing an automated semantic conflict detection system to evaluate draft ordinances from Davao City against a dual statutory search space. This collection encompasses national statutory enactments compiled to the extent accessible from public legal repositories, paired alongside the historical archive of enacted Davao City municipal ordinances collected to the fullest extent possible from local legislative records. Operationally, the system audits draft measures across two complementary dimensions: vertical statutory preemption, ensuring local provisions do not contradict superior national statutes under the \textit{Magtajas} doctrine and Republic Act No.~7160; and horizontal consistency, verifying that proposed clauses do not contradict, duplicate, or inadvertently cause implied repeals of active municipal ordinances.

To establish a feasible engineering boundary, several explicit delimitations govern this research, serving as baseline operational constraints while defining clear avenues for future work:
\begin{description}
    \item[Statutory Text Boundary.] The knowledge base is strictly bounded to codified national statutory instruments and enacted municipal ordinances. To keep retrieval indexes manageable and avoid ambiguous litigation records, judicial case law from Supreme Court decisions and departmental administrative circulars are excluded from the primary retrieval index. Synthesizing sprawling judicial case law alongside statutory enactments remains an open research frontier recommended for future large-scale legal informatics projects.

    \item[Municipal Archive Gaps.] The historical municipal archive reflects the practical realities of local record-keeping. While recent legislative terms benefit from digitized municipal database records, earlier decades rely primarily on curated landmark ordinances due to physical archive limitations. The system evaluates drafts strictly against the collection assembled through public portals and institutional coordination with the Sangguniang Panlungsod, leaving the OCR recovery of older paper records as an ongoing digitization task.

    \item[Linguistic Scope.] The system operates exclusively on English-language texts, following Section~20, Chapter~5, Book~I of Executive Order No.~292, which establishes that the English text controls in legal interpretation. While local drafts occasionally include regional terms or administrative Taglish, filtering the analysis to English avoids semantic drift and translation distortion. Extending conflict detection to multi-lingual and regional dialect reasoning is designated as a recommendation for future research.

    \item[Single-Premise Asymmetric Pairing.] The inference module evaluates proposed ordinance sections against individual candidate provisions using asymmetric span-level pairing ($1\text{-vs-}N$). The pipeline deliberately excludes multi-hop reasoning across multiple disparate enactments simultaneously. Identifying intricate legal contradictions that emerge only when synthesizing two or more separate national statutes exceeds the computational capacity of standard workstation hardware and is recommended for future exploration.

    \item[Bounded Open-Weight Model Envelope.] The evaluation is intentionally restricted to open-weight transformer encoder architectures capable of executing within the computational envelope of standard Local Government Unit (LGU) desktop computers (sub-600M parameters, sub-2GB inference VRAM). Massive autoregressive large language models (LLMs) and proprietary cloud APIs are excluded to ensure privacy, zero recurring cloud costs, and local offline deployment. Extending the pipeline to multi-billion parameter foundation models is reserved for future enterprise-grade studies.

    \item[Analytical Diagnostic Role.] The system functions strictly as a diagnostic decision-support tool providing white-box explainability, such as confidence metrics and attention heatmaps. The framework does not rewrite, generate, or alter legislative text, preserving human legal judgment in the hands of legislative researchers.

    \item[Direct Normative Scope vs. Longitudinal Statutory Construction.] The operational scope of the natural language inference engine is strictly delimited to direct normative and modal contradictions arising from explicit textual collisions (such as local prohibitions of federally authorized activities, unauthorized procedural preconditions, or explicit statutory penalty ceiling infractions). The system deliberately does not model longitudinal legislative histories, historical tax rate compounding, or multi-year financial calculations across unprovided statutory baselines, as observed in complex judicial tax disputes~\cite{scp2017mindanaoshopping}. Modeling such multi-hop temporal statutory construction and arithmetic accounting exceeds the ex-ante text-matching role of cross-encoder architectures and is preserved as an avenue for specialized fiscal analytics.
\end{description}

\section{Definition of Terms}
\begin{description}
    \item[Coarse-to-Fine Architecture.] A two-step computer process. First, the system quickly searches through thousands of documents to find "candidates" (coarse), then it carefully analyzes those specific candidates for logical conflicts (fine).

    \item[Contradiction.] A relationship in NLI where the second sentence is false if the first sentence is true, indicating a logical conflict.
    
    \item[Davao City Ordinances.] These are written local laws created by the Davao City government that are intended to be permanent and enforceable.
    
    \item[Entailment.] A relationship in NLI where the first sentence (premise) proves that the second sentence (hypothesis) is true.
    
    \item[Ex-Ante.] Derived from Latin meaning "before the event". Operationally, this refers to the evaluation or checking of a proposed city ordinance before it is officially signed into law.
    
    \item[Information Retrieval (IR).] The process of searching for and finding relevant documents (like specific laws) within a massive collection of text.
    
    \item[Local Government Code of 1991 (RA 7160).] The national law that gives local governments the power to create their own ordinances for the "General Welfare" of their citizens.
    
    \item[Magtajas v. Pryce Doctrine.] A legal rule from the Supreme Court stating that local ordinances are inferior to national laws and cannot go against any Act of Congress.
    
    \item[Natural Language Inference (NLI).] A task where the computer compares two sentences—a "premise" and a "hypothesis"—to see if they agree or disagree with each other.
    
    \item[Neutral.] A relationship in NLI where the two sentences are unrelated and do not affect each other's truth.
    
    \item[OCR Noise.] Text errors (like misspellings or broken words) that happen when physical paper documents are scanned into digital files.
    
    \item[Sangguniang Panlungsod (SP).] The local legislative body or "City Council" of Davao City, responsible for creating and approving local laws.
    
    \item[Semantic Conflict.] A logical clash where two different laws give opposite instructions. For example, a conflict occurs if one law permits an activity while another strictly prohibits it.
    
\end{description}
